% ---
% Reemplaza líneas 2722–2768 de Docs/Technical_Report.tex.
% Decisiones que lo respaldan: D-004, D-013, D-018, D-024, D-025, D-026,
% H-024, H-026, H-031.
% ---

\subsection{Módulo 4: Selección y escalado}

El módulo de selección y escalado prepara el vector de características para el circuito cuántico.
Debe reducir la dimensión al número de qubits, $k = 12$, y convertir cada característica en un
ángulo de rotación. Todos sus parámetros se ajustan solo con el conjunto de entrenamiento.

\textbf{Selección.} Las características se ordenan por el estadístico $F$ del ANOVA entre clases y
se recorren en ese orden: una característica se acepta solo si su correlación absoluta de Pearson
con todas las ya aceptadas es como máximo 0.95, hasta reunir doce. La restricción elimina
duplicados exactos que PyRadiomics produce en 2D, como \texttt{Energy} y \texttt{TotalEnergy}, y
variables casi idénticas, que harían que varios qubits codificaran el mismo valor. Sin ella, la
selección en masas no incluía ninguna característica de textura. Con ella, las doce características
de masas combinan intensidad, forma y textura (4, 5 y 3), y las de calcificaciones también (2, 3 y 7).

\textbf{Escalado angular.} Cada característica seleccionada, con valor radiómico $r_i$, se transforma
mediante su función de distribución empírica en entrenamiento, $\hat F_i$:
\begin{equation}
  x_i = \pi\,\hat F_i(r_i) \in [0,\pi].
  \label{eq:quantile_scaling}
\end{equation}
Es la transformación integral de probabilidad: cada ángulo se distribuye de manera uniforme en
$[0,\pi]$, por sesgada que sea la característica. Conserva el orden de las lesiones y descarta las
distancias dentro de cada característica. Con un escalado min-max, las características de cola
pesada, como la energía o la superficie, quedaban casi constantes: en calcificaciones, la energía de
la mitad de las lesiones caía a menos de 0.05 rad de cero, y su qubit recibía casi la misma rotación
para todas. El rango $[0,\pi]$ se debe a que el
\textit{feature map} aplica una fase $2x_i$: con $x_i\in[0,\pi]$, esa fase recorre la circunferencia
una sola vez, sin que dos valores distintos de la característica produzcan la misma rotación.

\textbf{Partición para la validación cruzada.} El conjunto de entrenamiento se divide en cinco
\textit{folds} estratificados por clase y agrupados por paciente, de modo que todas las lesiones de
un paciente quedan en el mismo \textit{fold}. La semilla de la partición se elige entre 200
candidatas con un criterio que solo mira las etiquetas: la menor desviación de la proporción de
malignas entre \textit{folds}. Los \textit{folds} se fijan aquí y los usan todos los módulos
posteriores.

El vector $\mathbf{x}\in[0,\pi]^{12}$ resultante es el mismo para las cinco condiciones de M5, de
forma que lo único que cambia entre ellas es la transformación.
